\subsection{What the shaped-charge example does and does not imply}
Proposition~\ref{prop:pathwise-interior} establishes a logical separation for a specified nonmonotone charge. It does not establish that pathwise randomization typically saves engineering or certification cost. The following result gives a sharp institutional boundary for that example rather than assigning its specially shaped charge an unsupported operational interpretation.
\begin{proposition}[One branch with a monotone charge]\label{prop:r43-monotone}
For continuously chosen levels and a single branch, concave terminal reward, convex intermediate cost, and a nonnegative nondecreasing charge per selected level, a deterministic one-level policy weakly dominates every randomized policy at each feasible promise. This holds for every alphabet budget and every realization tolerance, including the full pathwise institution with draw-specific intermediate service.
\end{proposition}
\begin{proof}
For an arbitrary feasible policy write $\mu=\E C$ and $\bar y=\E Y=B-\mu$. The deterministic level $\mu$ with intermediate service $\bar y$ attains total service $B\leq b$, and Jensen's inequalities improve both operating terms. The old book contains a largest used level $c_{\max}\geq\mu$. Nonnegativity and monotonicity imply that its total charge is at least $\rho(c_{\max})\geq\rho(\mu)$. Thus the replacement is feasible, uses one level, and does not lower net payoff.
\end{proof}
Consequently the one-branch separation in Proposition~\ref{prop:pathwise-interior} necessarily uses a charge outside this monotone class. The result covers zero charges, nondecreasing affine or convex charges, and a fixed opening fee plus a nondecreasing surcharge. It does not claim deterministic collapse for arbitrary heterogeneous multi-branch instances; the common alphabet couples their means. The original example and its proof remain intact as a mathematical possibility, not evidence of a generic operational advantage.
